% Every architecture figure and the model tables, at Springer's single-column
% text width (31 pc, about 13.1 cm), with draft captions.
% Build the figures first, then: pdflatex preview.tex (twice, for the numbers)
% In the paper: \includegraphics[width=\textwidth]{figures/arch/fig_f02.pdf}
\documentclass[10pt]{article}
\usepackage[a4paper, textwidth=31pc, textheight=194mm]{geometry}
\usepackage[T1]{fontenc}
\usepackage{newtxtext}
\usepackage{newtxmath}
\usepackage{graphicx}
\usepackage{booktabs}
\usepackage[labelfont=bf, font=small]{caption}
\pagestyle{empty}

\newcommand{\fig}[3]{%
  \begin{figure}[p]
    \centering
    \includegraphics[width=\linewidth, height=0.82\textheight, keepaspectratio]{#1}
    \caption{#2}\label{#3}
  \end{figure}
  \clearpage}

\begin{document}

\fig{fig_pipeline.pdf}{Preprocessing. Every clip becomes 15 frames of
224\,$\times$\,224 px centred on the striker. Steps 1--6 run once per clip;
only step 2 has learned weights, the dataset authors' YOLO11x-seg detector,
kept frozen. The example is a validation Cover Drive (vid533\_0): 28 broadcast
frames, 24 after removing repeats, 15 in the camera shot with the most
detections. The boxes are the ones the pipeline produced on frame 8.}{fig:pipeline}

\fig{fig_frame_models.pdf}{Frame models. One image backbone turns each of the
15 frames into a vector, a temporal head reads the 15 vectors, and a classifier
scores the 15 shots. The backbones are listed in Table~\ref{tab:image_backbones}
and the heads in Table~\ref{tab:temporal_heads}.}{fig:frame_models}

\fig{fig_video_models.pdf}{Video models. (a)~The clip is resampled to the
input its Kinetics-400 weights expect, and the backbone reads space and time
together (Table~\ref{tab:video_backbones}). (b)~R3D-18 reads motion instead of
colour: optical flow from RAFT-small (W01) or frame differences (W05). The
tiles are the real flow fields and differences of the example
clip.}{fig:video_models}

\fig{fig_pose_tracks.pdf}{The batter's body and the bat instead of pixels.
(a)~MediaPipe's 33 joints per frame, read by a GRU (E01) or a graph network
whose edges are the bones (E02). (b)~E04: each pose frame takes what it needs
from the RGB frames through cross-attention. (c)~W02--W04: the bat and ball
tracks from the dataset authors' detector. Every extractor is frozen; only the
heads train. The skeletons and bat boxes are the ones the models
read.}{fig:pose}

\fig{fig_f02.pdf}{F02, the multi-scale hybrid. (a)~Each frame passes through
the same multi-scale EfficientNetV2-S and a linear layer; a learned position is
added; two transformer encoder layers mix the 15 frames, and their mean goes
to the classifier. (b)~One encoder layer. (c)~Inside the frame encoder: the
maps after stages 3 and 5 are normalised, pooled to the last 7\,$\times$\,7
grid and stacked with it. 40.4\,M weights, half of them new. F01 is the same
encoder under GRU-128.}{fig:f02}

\fig{fig_ensemble.pdf}{The final system, chosen on the validation set before
the test set was read. The three runs of C06 and of F02 score the clip as
filmed and mirrored; each model's six answers are averaged, and the two models
vote with equal weight.}{fig:ensemble}

% The options behind the architecture figures, with the experiments that use
% them. Needs \usepackage{booktabs}. \input this file where the tables belong,
% or paste the tables in: Springer wants a single .tex file at submission.

\begin{table}[ht]
\centering
\caption{Image backbones of the frame models (Fig.~\ref{fig:frame_models}),
and how each one was used. A frozen backbone runs once and its frame vectors
are stored; a fine-tuned one trains with the head.}
\label{tab:image_backbones}
\begin{tabular}{@{}lll@{}}
\toprule
Backbone & Frozen & Fine-tuned end to end \\
\midrule
\multicolumn{3}{@{}l}{\textit{CNNs, pretrained on ImageNet}} \\
EfficientNetV2-S  & A01--A14, D01--D06, P01 & B01, F01--F05, R01--R03 \\
EfficientNetV2-B3 & -- & B02 \\
Xception          & -- & B03 \\
ConvNeXt-T        & P02 & B04 \\
MobileNetV3-L     & -- & B05 \\
ResNet-50         & -- & B06 \\
RegNetY-16GF      & -- & B07 \\
\multicolumn{3}{@{}l}{\textit{Vision transformers}} \\
Swin-T, ImageNet  & -- & B08 \\
CLIP ViT-B/16$^{a}$ & P05, V09 & -- \\
DINOv2-B$^{b}$ & V10 & -- \\
\bottomrule
\end{tabular}
\par\vspace{2pt}{\footnotesize $^{a}$ trained on 400\,M image--text pairs;\quad
$^{b}$ self-supervised, no labels.}
\end{table}

\begin{table}[ht]
\centering
\caption{Temporal heads of the frame models, grouped by how they treat the
order of the 15 frames. Each is followed by the same classifier: BatchNorm,
FC 1,024 with ReLU, FC 15.}
\label{tab:temporal_heads}
\begin{tabular}{@{}ll@{}}
\toprule
Temporal head & Experiments \\
\midrule
\multicolumn{2}{@{}l}{\textit{Recurrent: read the frames in order}} \\
GRU-128, last state & A03, B01--B08, R01--R03 \\
GRU-64 / BiGRU-128 & A04 / A05 \\
LSTM-64 / BiLSTM-64 & A06 / A07 \\
GRU-128 after a LayerNorm & A13 \\
GRU-128 on averaged feature maps & A14 \\
GRU-128 with attention over its steps & F05 \\
\multicolumn{2}{@{}l}{\textit{Attention: each frame attends to all 15}} \\
Transformer, 2 layers, 4 heads & A08, F02 \\
Transformer, 4 layers, 8 heads & A09 \\
\multicolumn{2}{@{}l}{\textit{Convolution over time}} \\
Dilated temporal convolution, rates 1, 2, 4 & A10 \\
\multicolumn{2}{@{}l}{\textit{Pooling: frame order ignored}} \\
Mean / max over frames & A01 / A02 \\
Attention pooling & A11 \\
NetVLAD, 8 learned centres & A12 \\
\multicolumn{2}{@{}l}{\textit{No temporal head: the frame vectors averaged}} \\
One linear layer (a probe) & P01, P02, P05 \\
MLP, 1,024 hidden & V09, V10 \\
Logistic regression, RBF SVM, random forest, & D01--D06 \\
\quad XGBoost, 10-NN, MLP & \\
\bottomrule
\end{tabular}
\end{table}

\begin{table}[ht]
\centering
\caption{Video backbones (Fig.~\ref{fig:video_models}a), the input each one
takes, and how it was used. Our 15 frames become 16 by repeating the middle
one and 8 by keeping every other one.}
\label{tab:video_backbones}
\begin{tabular}{@{}lccccc@{}}
\toprule
Backbone & Frames & Size (px) & Fine-tuned, & Frozen, & Frozen, \\
 & & & new linear & MLP & one linear \\
\midrule
\multicolumn{6}{@{}l}{\textit{3D CNNs, Kinetics-400}} \\
R3D-18 & 15 & 112 & C02 & V01 & P03 \\
MC3-18 & 15 & 112 & -- & V02 & -- \\
R(2+1)D-18 & 15 & 112 & C01 & V03 & -- \\
S3D & 15 & 224 & -- & V04 & -- \\
I3D-R50 & 8 & 224 & C03 & -- & -- \\
SlowFast-R50 & 32 + 8$^{a}$ & 224 & C04 & -- & -- \\
X3D-M & 16 & 224 & C05 & -- & -- \\
\multicolumn{6}{@{}l}{\textit{Video transformers, Kinetics-400}} \\
MViTv2-S & 16 & 224 & C06 & V06 & -- \\
Swin3D-T & 15 & 224 & C07 & V05 & -- \\
VideoMAE-B & 16 & 224 & C08 & V07 & P04 \\
TimeSformer-B & 8 & 224 & -- & V08 & -- \\
\bottomrule
\end{tabular}
\par\vspace{2pt}{\footnotesize $^{a}$ SlowFast reads the clip twice: 32 frames on its fast path, 8 on its slow one.}
\end{table}


\end{document}
